\hypertarget{appendix-7bc181935c42}{}
\section{Derivative as a local linear map}\label{appendix-7bc181935c42-derivative-as-a-local-linear-map}

\textbf{Curriculum depth:} Developed · \textbf{First seen:} Phase 9 · \textbf{Latest development:} Phase 9

\subsection{Definition and mechanism}\label{appendix-7bc181935c42-definition-and-mechanism}

The derivative measures first-order change. In one variable it is the limit of the difference quotient; in several variables it is the linear map that best predicts small output changes near a point. For a scalar function, this local model is the tangent line. For a vector-valued function, it is a matrix acting on an input perturbation. The derivative is local: its approximation error becomes small relative to the perturbation as the step shrinks, but the linear model need not remain accurate over a large region. This view unifies instantaneous velocity, gradients, Jacobians, and optimization directions. It also gives a practical check: compare a small finite perturbation with the derivative's predicted change, while remembering that floating-point rounding and curvature limit how small or large the perturbation should be.

\subsection{Prerequisites}\label{appendix-7bc181935c42-prerequisites}

\begin{itemize}
\tightlist
\item
  No earlier concept node is required in this planned path.
\end{itemize}

\subsection{Used by}\label{appendix-7bc181935c42-used-by}

\begin{itemize}
\tightlist
\item
  \href{../03-Visual-Computing-I/main.pdf\#nameddest=appendix-40df0855b4a7}{Automatic differentiation and backpropagation (03-Visual-Computing-I, appendix-40df0855b4a7)}
\item
  \hyperref[appendix-1a50e72435c2]{Definite integrals and the fundamental theorem}
\item
  \hyperref[appendix-e2d476574aab]{Directional and total derivatives}
\item
  \hyperref[appendix-e2d7b4f2161d]{Finite differences and step-size balance}
\item
  \href{../01-Foundations-I/main.pdf\#nameddest=appendix-d8a7ddb675eb}{Jacobians, pushforwards, and pullbacks (01-Foundations-I, appendix-d8a7ddb675eb)}
\item
  \href{../06-Engine-II/main.pdf\#nameddest=appendix-677b59746ae0}{Partial derivatives and gradients (06-Engine-II, appendix-677b59746ae0)}
\item
  \hyperref[appendix-32f39e9a3738]{Taylor models and approximation error}
\end{itemize}

\subsection{Full teaching explanation}\label{appendix-7bc181935c42-full-teaching-explanation}

\begin{itemize}
\tightlist
\item
  \href{../01-Foundations-I/main.pdf\#nameddest=01-foundations-i-ch020-functions-limits-continuity-and-scalar-change}{Phase 9 teaching section (01-Foundations-I, 01-foundations-i-ch020-functions-limits-continuity-and-scalar-change)}
\item
  \href{../01-Foundations-I/main.pdf\#nameddest=appendix-adf48f5dc5c1}{Visual Recall/09 - Calculus and Optimization - Visual Recall.md (01-Foundations-I, appendix-adf48f5dc5c1)} · visual recall
\end{itemize}

\subsection{Source coverage}\label{appendix-7bc181935c42-source-coverage}

Supplied lesson parts: 805 (T04 · M049), 806 (T04 · M049), 807 (T04 · M049).
